% Source: https://artofproblemsolving.com/community/c3h2277308_logic_problem
% Stable identity: aops-topic-2277308-q01
The exchange condition means each person has at most $4$ pennies, at most $1$ nickel ($2$ nickels make a dime), at most $1$ quarter ($2$ quarters make a half-dollar) and at most $4$ dimes ($5$ dimes make a half-dollar). With these limits, the number of pennies is forced by the remainder of the amount modulo $5$.

\textbf{Sam, $49$ cents.} He has $4$ pennies, so $45$ cents in nickels, dimes and quarters. With a quarter, $20$ cents remain, which with at most one nickel can only be $2$ dimes: $7$ coins in all. Without a quarter, $45$ cents must be $4$ dimes and a nickel: $9$ coins. So Sam has at least $7$ coins (a quarter, two dimes and four pennies).

\textbf{Sally, $43$ cents.} She has $3$ pennies and $40$ cents in nickels, dimes and quarters: either a quarter, a dime and a nickel ($3$ coins) or $4$ dimes ($4$ coins). So Sally has at least $3+3=6$ coins.

Together they have at least $7+6=\boxed{13}$ coins.
